\chapter{L4 -- Soluzione dei sistemi lineari: caso discreto generale, caso stazionario, successione di Peano--Baker ed esponenziale di matrice}

\section{Collegamento tra forma esplicita e forma implicita}

Nella lezione precedente abbiamo ottenuto le relazioni che collegano le quantità ``globali'' $\Phi$ e $H$ alle quantità ``locali'' $A$ e $B$.

\subsection{Tempo continuo}

Per il tempo continuo valgono:
\begin{equation}
A(t)=\left.\frac{\partial \Phi(\xi,t)}{\partial \xi}\right|_{\xi=t},
\qquad
B(t)=H(t,t).
\label{eq:L4_def_A_B_TC}
\end{equation}

Quindi la rappresentazione implicita locale del sistema è
\begin{equation}
\dot x(t)=A(t)x(t)+B(t)u(t),
\qquad
y(t)=C(t)x(t)+D(t)u(t).
\label{eq:L4_TC_implicita}
\end{equation}

\subsection{Tempo discreto}

Nel tempo discreto, invece, si ha:
\begin{equation}
A(j)=\Phi(j+1,j),
\qquad
B(j)=H(j+1,j),
\label{eq:L4_def_A_B_TD}
\end{equation}
e dunque la forma implicita diventa
\begin{equation}
x(k+1)=A(k)x(k)+B(k)u(k),
\qquad
y(k)=C(k)x(k)+D(k)u(k).
\label{eq:L4_TD_implicita}
\end{equation}

L'obiettivo di questa lezione è capire come, a partire da \eqref{eq:L4_TD_implicita} e \eqref{eq:L4_TC_implicita}, si ricostruiscono le forme esplicite e come si presentano le matrici $\Phi$ e $H$.

\section{Sistemi lineari discreti generici}

Consideriamo il sistema discreto non necessariamente stazionario
\begin{equation}
x(k+1)=A(k)x(k)+B(k)u(k),
\qquad
x(K_0)\ \text{assegnato}.
\label{eq:L4_sys_TD}
\end{equation}

\subsection{Sviluppo iterativo dei primi passi}

Iteriamo la dinamica a partire da $K_0$.

Per $k=K_0$:
\[
x(K_0+1)=A(K_0)x(K_0)+B(K_0)u(K_0).
\]

Per $k=K_0+1$:
\begin{align*}
x(K_0+2)
&=A(K_0+1)x(K_0+1)+B(K_0+1)u(K_0+1)\\
&=A(K_0+1)\bigl[A(K_0)x(K_0)+B(K_0)u(K_0)\bigr]+B(K_0+1)u(K_0+1)\\
&=A(K_0+1)A(K_0)x(K_0)+A(K_0+1)B(K_0)u(K_0)+B(K_0+1)u(K_0+1).
\end{align*}

Per $k=K_0+2$:
\begin{align*}
x(K_0+3)
&=A(K_0+2)x(K_0+2)+B(K_0+2)u(K_0+2)\\
&=A(K_0+2)A(K_0+1)A(K_0)x(K_0)\\
&\quad +A(K_0+2)A(K_0+1)B(K_0)u(K_0)\\
&\quad +A(K_0+2)B(K_0+1)u(K_0+1)\\
&\quad +B(K_0+2)u(K_0+2).
\end{align*}

A questo punto il pattern è chiaro:
\begin{itemize}
    \item compare un termine libero, cioè la propagazione dello stato iniziale;
    \item compare una somma di termini forzati, uno per ogni ingresso passato, ciascuno trasportato fino all'istante finale.
\end{itemize}

\subsection{Forma esplicita generale}

Introduciamo la convenzione di prodotto ordinato
\[
\prod_{i=K_0}^{K-1}A(i):=A(K-1)A(K-2)\cdots A(K_0),
\]
cioè gli indici crescono da destra verso sinistra, come richiesto dalla composizione delle matrici di stato.

Con questa convenzione si ottiene:
\begin{equation}
x(K)=\Phi(K,K_0)x(K_0)+\sum_{j=K_0}^{K-1}H(K,j)u(j),
\label{eq:L4_explicit_TD}
\end{equation}
dove
\begin{equation}
\Phi(K,K_0)=\prod_{i=K_0}^{K-1}A(i),
\label{eq:L4_Phi_TD_generic}
\end{equation}
e
\begin{equation}
H(K,j)=\left(\prod_{i=j+1}^{K-1}A(i)\right)B(j).
\label{eq:L4_H_TD_generic}
\end{equation}

\paragraph{Convenzione importante.}
Se il prodotto è vuoto, esso vale l'identità:
\[
\prod_{i=K}^{K-1}A(i)=I.
\]
Per esempio, per $j=K-1$ si ha
\[
H(K,K-1)=B(K-1).
\]

\paragraph{Interpretazione di $H(K,j)$.}
Il termine $H(K,j)$ dice come un ingresso applicato all'istante $j$ viene trasportato fino all'istante finale $K$.  
In altre parole:
\begin{itemize}
    \item $B(j)$ trasforma l'ingresso in variazione immediata dello stato;
    \item il prodotto di matrici $A$ successive trasporta quella variazione fino al tempo finale.
\end{itemize}

\begin{figure}[H]
\centering
\begin{tikzpicture}[>=Latex, node distance=1.45cm]
    \node[draw, rounded corners, minimum width=1.2cm, minimum height=0.9cm] (x0) {$x(K_0)$};
    \node[draw, rounded corners, right=of x0, minimum width=1.2cm, minimum height=0.9cm] (x1) {$x(K_0+1)$};
    \node[draw, rounded corners, right=of x1, minimum width=1.2cm, minimum height=0.9cm] (x2) {$x(K_0+2)$};
    \node[right=0.7cm of x2] (dots) {$\cdots$};
    \node[draw, rounded corners, right=0.7cm of dots, minimum width=1.2cm, minimum height=0.9cm] (xk) {$x(K)$};

    \draw[->, thick] (x0) -- (x1) node[midway, above] {$A(K_0)$};
    \draw[->, thick] (x1) -- (x2) node[midway, above] {$A(K_0+1)$};
    \draw[->, thick] (x2) -- (dots);
    \draw[->, thick] (dots) -- (xk) node[midway, above] {$A(K-1)$};

    \draw[->, thick] ($(x0)+(0,-1.4)$) -- (x1.south) node[midway, left] {$B(K_0)u(K_0)$};
    \draw[->, thick] ($(x1)+(0,-1.4)$) -- (x2.south) node[midway, left] {$B(K_0+1)u(K_0+1)$};
\end{tikzpicture}
\caption{Interpretazione della dinamica discreta: lo stato iniziale si propaga tramite i prodotti di $A(k)$, mentre ciascun ingresso entra attraverso $B(k)$ e poi viene trasportato in avanti.}
\end{figure}

\section{Caso discreto stazionario}

Consideriamo ora il caso in cui il sistema sia tempo-invariante:
\[
A(k)=A,\qquad B(k)=B,\qquad C(k)=C,\qquad D(k)=D
\qquad \forall k.
\]

Allora le formule precedenti si semplificano molto.

\subsection{Matrice di transizione}

Da \eqref{eq:L4_Phi_TD_generic} segue
\begin{equation}
\Phi(K,K_0)=A^{K-K_0}.
\label{eq:L4_Phi_TD_LTI}
\end{equation}

\subsection{Nucleo forzato}

Da \eqref{eq:L4_H_TD_generic} segue
\begin{equation}
H(K,j)=A^{K-j-1}B.
\label{eq:L4_H_TD_LTI}
\end{equation}

\subsection{Soluzione esplicita}

Quindi la soluzione del sistema discreto stazionario è
\begin{equation}
x(K)=A^{K-K_0}x(K_0)+\sum_{j=K_0}^{K-1}A^{K-j-1}Bu(j).
\label{eq:L4_sol_TD_LTI}
\end{equation}

L'uscita vale
\begin{equation}
y(K)=Cx(K)+Du(K).
\label{eq:L4_out_TD_LTI}
\end{equation}

\paragraph{Messaggio chiave.}
Nel caso discreto stazionario, tutto il comportamento del sistema dipende dallo studio delle potenze di $A$.  
Per questo motivo il corso dedicherà grande attenzione alle proprietà spettrali di $A$.

\section{Caso continuo: dinamica libera e matrice di transizione}

Passiamo ora al tempo continuo:
\begin{equation}
\dot x(t)=A(t)x(t)+B(t)u(t).
\label{eq:L4_sys_TC}
\end{equation}

Per capire come è fatta la matrice di transizione, si parte dall'evoluzione libera, cioè dal caso
\[
u(t)=0.
\]

In questo caso la dinamica diventa
\begin{equation}
\dot x(t)=A(t)x(t),
\qquad x(t_0)=x_0.
\label{eq:L4_free_TC}
\end{equation}

Poiché l'evoluzione libera è descritta da $\Phi(t,t_0)$, deve valere
\[
x(t)=\Phi(t,t_0)x(t_0).
\]

Derivando rispetto a $t$:
\[
\dot x(t)=\frac{\partial \Phi(t,t_0)}{\partial t}x(t_0).
\]

Ma dalla dinamica libera \eqref{eq:L4_free_TC} si ha anche
\[
\dot x(t)=A(t)x(t)=A(t)\Phi(t,t_0)x(t_0).
\]

Poiché la relazione vale per ogni $x(t_0)$, si ottiene l'\emph{equazione differenziale matriciale}:
\begin{equation}
\boxed{
\frac{\partial \Phi(t,t_0)}{\partial t}=A(t)\Phi(t,t_0),
\qquad
\Phi(t_0,t_0)=I.
}
\label{eq:L4_matrix_ode}
\end{equation}

\paragraph{Questa è la caratterizzazione fondamentale di $\Phi$.}
Essa dice che trovare la matrice di transizione equivale a risolvere un problema di Cauchy matriciale.

\section{Successione di Peano}

Per risolvere \eqref{eq:L4_matrix_ode} in modo costruttivo si introduce la \emph{successione di Peano}.

Si cerca $\Phi(t,t_0)$ sotto forma di serie
\begin{equation}
\Phi(t,t_0)=\sum_{k=0}^{\infty} S_k(t,t_0),
\label{eq:L4_peano_series}
\end{equation}
dove i termini sono definiti ricorsivamente da
\begin{equation}
S_0(t,t_0)=I,
\label{eq:L4_S0}
\end{equation}
\begin{equation}
S_1(t,t_0)=\int_{t_0}^{t}A(\tau)S_0(\tau,t_0)\,d\tau
=
\int_{t_0}^{t}A(\tau)\,d\tau,
\label{eq:L4_S1}
\end{equation}
e, più in generale,
\begin{equation}
S_{k+1}(t,t_0)=\int_{t_0}^{t}A(\tau)S_k(\tau,t_0)\,d\tau.
\label{eq:L4_Sk_rec}
\end{equation}

\paragraph{Interpretazione.}
Ogni termine della successione rappresenta un ordine successivo di ``accumulo'' dell'effetto della matrice $A(\tau)$ lungo l'intervallo di integrazione.

\subsection{Primi termini}

Scriviamo i primi termini per capire la struttura:
\[
S_0(t,t_0)=I,
\]
\[
S_1(t,t_0)=\int_{t_0}^{t}A(\tau_1)\,d\tau_1,
\]
\[
S_2(t,t_0)=\int_{t_0}^{t}A(\tau_1)\left(\int_{t_0}^{\tau_1}A(\tau_2)\,d\tau_2\right)d\tau_1,
\]
e così via.

In modo informale, la successione di Peano è l'analogo matriciale della serie iterativa che si ottiene riscrivendo il problema differenziale come equazione integrale.

\section{Stima dei termini della successione di Peano}

Per dimostrare che la serie \eqref{eq:L4_peano_series} converge, serve una stima uniforme sui termini $S_k$.

Sia
\[
M=\sup_{\tau\in[t_0,t]}\norm{A(\tau)}.
\]
Supponiamo che $A(\tau)$ sia limitata sull'intervallo considerato. Allora vale il seguente risultato.

\begin{proposizione}
Per ogni $k\ge 0$,
\begin{equation}
\norm{S_k(t,t_0)}\le \frac{M^k(t-t_0)^k}{k!}.
\label{eq:L4_bound_Sk}
\end{equation}
\end{proposizione}

\begin{proof}
Procediamo per induzione.

Per $k=0$:
\[
S_0(t,t_0)=I,
\]
e la stima vale banalmente, perché
\[
\norm{S_0(t,t_0)}=\norm{I}\le 1.
\]
Nel seguito si assume una norma matriciale compatibile e, se necessario, si assorbe il valore di $\norm{I}$ nella costante.

Per $k=1$:
\[
S_1(t,t_0)=\int_{t_0}^{t}A(\tau)\,d\tau.
\]
Usando disuguaglianza triangolare e definizione di $M$:
\[
\norm{S_1(t,t_0)}
\le
\int_{t_0}^{t}\norm{A(\tau)}\,d\tau
\le
\int_{t_0}^{t}M\,d\tau
=
M(t-t_0).
\]

Supponiamo ora vera la stima per un certo $k$ e dimostriamola per $k+1$:
\[
S_{k+1}(t,t_0)=\int_{t_0}^{t}A(\tau)S_k(\tau,t_0)\,d\tau.
\]
Allora
\begin{align*}
\norm{S_{k+1}(t,t_0)}
&\le
\int_{t_0}^{t}\norm{A(\tau)}\,\norm{S_k(\tau,t_0)}\,d\tau\\
&\le
\int_{t_0}^{t}
M\,
\frac{M^k(\tau-t_0)^k}{k!}
\,d\tau\\
&=
\frac{M^{k+1}}{k!}\int_{t_0}^{t}(\tau-t_0)^k\,d\tau\\
&=
\frac{M^{k+1}}{k!}\cdot \frac{(t-t_0)^{k+1}}{k+1}\\
&=
\frac{M^{k+1}(t-t_0)^{k+1}}{(k+1)!}.
\end{align*}
La tesi è dimostrata.
\end{proof}

\paragraph{Conseguenza.}
La serie di Peano è dominata dalla serie esponenziale
\[
\sum_{k=0}^{\infty}\frac{M^k(t-t_0)^k}{k!}=e^{M(t-t_0)},
\]
quindi converge assolutamente e uniformemente su ogni intervallo limitato.

\section{Verifica che la serie di Peano risolve il problema matriciale}

Una volta definita
\[
\Phi(t,t_0)=\sum_{k=0}^{\infty}S_k(t,t_0),
\]
bisogna verificare che soddisfi davvero \eqref{eq:L4_matrix_ode}.

\subsection{Condizione iniziale}

Valutiamo la serie per $t=t_0$:
\[
\Phi(t_0,t_0)=\sum_{k=0}^{\infty}S_k(t_0,t_0).
\]

Ora:
\[
S_0(t_0,t_0)=I,
\]
mentre per ogni $k\ge 1$ i termini sono integrali su intervalli degeneri, quindi
\[
S_k(t_0,t_0)=0.
\]

Ne segue
\begin{equation}
\Phi(t_0,t_0)=I.
\label{eq:L4_Phi_initial}
\end{equation}

\subsection{Equazione differenziale}

Dalla definizione ricorsiva:
\[
S_k(t,t_0)=\int_{t_0}^{t}A(\tau)S_{k-1}(\tau,t_0)\,d\tau,
\qquad k\ge 1.
\]

Derivando termine a termine:
\[
\frac{\partial S_k(t,t_0)}{\partial t}=A(t)S_{k-1}(t,t_0),\qquad k\ge 1.
\]

Allora
\begin{align*}
\frac{\partial \Phi(t,t_0)}{\partial t}
&=
\sum_{k=1}^{\infty}\frac{\partial S_k(t,t_0)}{\partial t}\\
&=
\sum_{k=1}^{\infty}A(t)S_{k-1}(t,t_0)\\
&=
A(t)\sum_{k=1}^{\infty}S_{k-1}(t,t_0)\\
&=
A(t)\sum_{j=0}^{\infty}S_j(t,t_0)\\
&=
A(t)\Phi(t,t_0).
\end{align*}

Quindi
\begin{equation}
\frac{\partial \Phi(t,t_0)}{\partial t}=A(t)\Phi(t,t_0),
\end{equation}
come volevamo dimostrare.

\section{Espressione di \texorpdfstring{$H(t,\tau)$}{H(t,tau)} nel continuo}

Nella lezione precedente era emersa la proprietà di composizione
\[
\Phi(t,\tau)H(\tau,t_0)=H(t,t_0).
\]

Scegliendo come istante iniziale il tempo $\tau$, si ottiene
\[
H(t,\tau)=\Phi(t,\tau)H(\tau,\tau).
\]

Ma per definizione
\[
H(\tau,\tau)=B(\tau).
\]

Quindi
\begin{equation}
H(t,\tau)=\Phi(t,\tau)B(\tau).
\label{eq:L4_H_TC_from_Phi}
\end{equation}

Sostituendo nella rappresentazione esplicita dello stato:
\begin{equation}
x(t)=\Phi(t,t_0)x(t_0)+\int_{t_0}^{t}\Phi(t,\tau)B(\tau)u(\tau)\,d\tau.
\label{eq:L4_explicit_TC_general}
\end{equation}

Questa è la formula generale della soluzione dei sistemi lineari a tempo continuo.

\section{Caso continuo stazionario}

Supponiamo ora che il sistema sia tempo-invariante:
\[
A(t)=A,\qquad B(t)=B,\qquad C(t)=C,\qquad D(t)=D.
\]

In questo caso la ricorsione di Peano si semplifica moltissimo.

\subsection{Termini della successione di Peano}

Poiché $A(\tau)=A$ è costante, i primi termini diventano:
\[
S_0(t,t_0)=I,
\]
\[
S_1(t,t_0)=\int_{t_0}^{t}A\,d\tau=A(t-t_0),
\]
\[
S_2(t,t_0)=\int_{t_0}^{t}A\,S_1(\tau,t_0)\,d\tau
=
\int_{t_0}^{t}A\cdot A(\tau-t_0)\,d\tau
=
A^2\frac{(t-t_0)^2}{2!}.
\]

Procedendo per induzione si ottiene
\begin{equation}
S_k(t,t_0)=\frac{A^k(t-t_0)^k}{k!}.
\label{eq:L4_Peano_TI_term}
\end{equation}

\subsection{Matrice di transizione}

Allora la serie di Peano diventa
\[
\Phi(t,t_0)=\sum_{k=0}^{\infty}\frac{A^k(t-t_0)^k}{k!}.
\]

Questa è precisamente la definizione di \emph{esponenziale di matrice}:
\begin{equation}
\boxed{
\Phi(t,t_0)=e^{A(t-t_0)}.
}
\label{eq:L4_Phi_exp_matrix}
\end{equation}

\paragraph{Osservazione fondamentale.}
L'esponenziale di matrice gioca, per i sistemi lineari continui, lo stesso ruolo che le potenze di $A$ giocano per i sistemi lineari discreti stazionari.

\subsection{Nucleo forzato}

Da \eqref{eq:L4_H_TC_from_Phi} segue
\[
H(t,\tau)=\Phi(t,\tau)B=e^{A(t-\tau)}B.
\]

\subsection{Soluzione completa}

La soluzione del sistema lineare continuo stazionario è quindi
\begin{equation}
x(t)=e^{A(t-t_0)}x(t_0)+\int_{t_0}^{t}e^{A(t-\tau)}Bu(\tau)\,d\tau.
\label{eq:L4_sol_TC_LTI}
\end{equation}

L'uscita vale
\begin{equation}
y(t)=Cx(t)+Du(t).
\label{eq:L4_out_TC_LTI}
\end{equation}

\begin{figure}[H]
\centering
\begin{tikzpicture}[>=Latex]
    \draw[->] (0,0) -- (9.8,0) node[right] {$\tau$};
    \draw[->] (0,0) -- (0,2.6) node[above] {};
    \draw[dashed] (1.5,0) -- (1.5,2.1) node[above] {$t_0$};
    \draw[dashed] (8.2,0) -- (8.2,2.1) node[above] {$t$};

    \draw[thick,blue] (2.3,0.6) -- (2.3,1.3);
    \draw[thick,blue] (4.2,0.6) -- (4.2,1.6);
    \draw[thick,blue] (6.1,0.6) -- (6.1,1.1);

    \node[blue] at (2.3,1.55) {$u(\tau_1)$};
    \node[blue] at (4.2,1.85) {$u(\tau_2)$};
    \node[blue] at (6.1,1.35) {$u(\tau_3)$};

    \draw[->, thick,red] (2.3,1.3) to[bend left=15] (8.0,2.2);
    \draw[->, thick,red] (4.2,1.6) to[bend left=10] (8.0,2.0);
    \draw[->, thick,red] (6.1,1.1) to[bend left=5]  (8.0,1.8);

    \node[red] at (8.8,2.25) {$x(t)$};
\end{tikzpicture}
\caption{Interpretazione della convoluzione nel continuo: ogni ingresso passato $u(\tau)$ contribuisce allo stato finale $x(t)$ attraverso il peso $\Phi(t,\tau)B$.}
\end{figure}

\section{Confronto finale tra continuo e discreto}

A questo punto il parallelo tra i due casi è completo.

\subsection{Tempo discreto stazionario}

\[
x(K)=A^{K-K_0}x(K_0)+\sum_{j=K_0}^{K-1}A^{K-j-1}Bu(j),
\]
\[
y(K)=Cx(K)+Du(K).
\]

\subsection{Tempo continuo stazionario}

\[
x(t)=e^{A(t-t_0)}x(t_0)+\int_{t_0}^{t}e^{A(t-\tau)}Bu(\tau)\,d\tau,
\]
\[
y(t)=Cx(t)+Du(t).
\]

\paragraph{Analogia concettuale.}
\begin{itemize}
    \item Nel discreto l'evoluzione libera è governata da $A^k$.
    \item Nel continuo l'evoluzione libera è governata da $e^{At}$.
    \item Nel discreto la risposta forzata è una somma.
    \item Nel continuo la risposta forzata è un integrale di convoluzione.
\end{itemize}

\section{Sintesi della lezione}

In questa lezione abbiamo visto che:

\begin{enumerate}
    \item nel tempo discreto, iterando la dinamica
    \[
    x(k+1)=A(k)x(k)+B(k)u(k),
    \]
    si ottengono esplicitamente
    \[
    \Phi(K,K_0)=\prod_{i=K_0}^{K-1}A(i),
    \qquad
    H(K,j)=\left(\prod_{i=j+1}^{K-1}A(i)\right)B(j);
    \]

    \item nel caso discreto stazionario queste formule si riducono a
    \[
    \Phi(K,K_0)=A^{K-K_0},
    \qquad
    H(K,j)=A^{K-j-1}B;
    \]

    \item nel tempo continuo la matrice di transizione è la soluzione del problema
    \[
    \dot \Phi(t,t_0)=A(t)\Phi(t,t_0),
    \qquad
    \Phi(t_0,t_0)=I;
    \]

    \item tale soluzione si costruisce tramite la successione di Peano;

    \item nel caso tempo-invariante la serie di Peano coincide con l'esponenziale di matrice:
    \[
    \Phi(t,t_0)=e^{A(t-t_0)};
    \]

    \item la soluzione completa del sistema continuo stazionario è
    \[
    x(t)=e^{A(t-t_0)}x(t_0)+\int_{t_0}^{t}e^{A(t-\tau)}Bu(\tau)\,d\tau.
    \]
\end{enumerate}

Questi risultati sono fondamentali: da qui in avanti tutto lo studio dei sistemi lineari ruoterà intorno a
\begin{itemize}
    \item potenze di matrici nel discreto;
    \item esponenziale di matrice nel continuo;
    \item struttura spettrale della matrice $A$.
\end{itemize}